\section{Introduction}
\label{sec:intro}

\IEEEPARstart{F}{alls} rank among the commonest causes of workplace death. In the United Kingdom, falls from height were the most frequent fatal accident to workers in 2025/26, about a quarter of all worker deaths \cite{hse2026fatal}; in the United States, falls, slips and trips caused 844 of the 5,070 fatal work injuries in 2024 \cite{bls2026cfoi}. An injured worker is often unable to summon help, and the outcome hinges on how swiftly someone notices. Wearable ``man-down'' alarms, which UK lone-working guidance recommends \cite{hse2020loneworkers}, mitigate this risk but depend on the worker wearing and charging a device. Fixed cameras already blanket many industrial floors, so detecting a person lying on the floor from video is a natural complement.

Two constraints govern such a system. It must run close to the cameras, since streaming every feed to a data centre consumes bandwidth, adds latency and moves footage of employees off site, which many sites will not tolerate; the detector must therefore fit on an inexpensive edge device serving several cameras. And it must separate a person who is down from postures that look alike in a single frame: a mechanic kneeling at a machine, a picker crouching at a low shelf, a worker sitting on the floor, someone lying under a vehicle. False alarms from such postures are why these systems end up ignored.

Classic fall detectors either learn a classifier from labelled videos \cite{mubashir2013survey,igual2013challenges} or apply hand-written rules to body keypoints \cite{cao2021openpose,maji2022yolopose}. Both excel on the datasets they were built for, yet both demand fresh labels or rules whenever the scene, camera angle or set of confusable postures shifts. Vision--language models (VLMs) offer an alternative: the task is stated in words, and the model answers without task-specific training \cite{radford2021clip,liu2023llava}. VLMs with under one billion parameters now run on edge hardware \cite{marafioti2025smolvlm,amini2025lfm2}, and integer quantization shrinks them further \cite{dettmers2022llmint8,frantar2023gptq,lin2024awq}. What is missing is evidence of where they cease to suffice. How much accuracy is forfeited in moving from 4\,B to 450\,M parameters, from 16-bit to 4-bit weights, or from full resolution to 256 pixels? Does a short clip beat a larger model? How many cameras can one device sustain?

We answer these questions for \pd{} detection in a zero-shot study: no model is fine-tuned, so the results describe what a practitioner obtains by downloading a model and writing a prompt. Our research questions are:
\begin{itemize}
  \item \textbf{RQ1} How do accuracy and false-alarm rate vary with model size, quantization, input resolution and number of frames per decision?
  \item \textbf{RQ2} What are the latency, peak memory and energy costs of each configuration on a consumer device, on its CPU alone and on its integrated GPU?
  \item \textbf{RQ3} Which configurations lie on the accuracy--latency and accuracy--memory Pareto frontiers?
  \item \textbf{RQ4} Can VLMs tell a person who is down from hard negatives such as kneeling, crouching, sitting on the floor and lying on furniture?
\end{itemize}

Our contributions are:
\begin{enumerate}
  \item A \pd{} benchmark of \BenchNKept{} frames and clips drawn from two public fall datasets \cite{kwolek2014urfall,auvinet2010mcfd}, with a written labelling policy, every frame checked by eye\iffullcheck\else\pending{after the author's full check}\fi{}, a blind agreement check, six hard-negative posture types, viewpoint annotations and sequence-disjoint development and test splits. We release labels and build scripts; the frames remain with their original distributors.
  \item A zero-shot evaluation of five open VLMs (450\,M--4\,B parameters) across quantization, resolution, temporal context and prompt wording, plus a keypoint-rule baseline that tests whether a VLM is needed at all.
  \item Latency, memory and energy measurements on the CPU and the integrated GPU of an Apple M4 Pro under a documented protocol, including two measurement pitfalls on unified-memory hardware that can distort reported memory by more than a factor of three.
  \item Accuracy--efficiency Pareto frontiers, a table of the camera streams each configuration sustains at one decision per second, and deployment recommendations for each device class.
\end{enumerate}

Section~\ref{sec:related} reviews related work. Section~\ref{sec:benchmark} describes the benchmark, Section~\ref{sec:methods} the models and baseline, and Section~\ref{sec:protocol} the evaluation protocol. Section~\ref{sec:results} reports results, Section~\ref{sec:discussion} discusses them with the study's limitations, and Section~\ref{sec:conclusion} concludes.
